\documentclass[a4paper]{article}

\usepackage[dutch]{babel}
\usepackage[pdftex]{graphicx}
\usepackage{amsfonts}
\usepackage{amsmath}
\usepackage{amssymb}
\usepackage{amsthm}
\usepackage{xcolor}
\usepackage{mdframed}
\usepackage{parskip}
\usepackage{array}
\usepackage{enumitem}
\usepackage{bm}
\usepackage{tikz}
\usetikzlibrary{arrows,automata,calc,fit,positioning}
\usepackage{chngcntr}
\counterwithin{figure}{section}

\begingroup
    \makeatletter
    \@for\theoremstyle:=definition,remark,plain\do{%
        \expandafter\g@addto@macro\csname th@\theoremstyle\endcsname{%
            \addtolength\thm@preskip\parskip
            }%
        }
\endgroup

\newcommand\addvmargin[1]{
  \node[fit=(current bounding box),inner ysep=#1,inner xsep=0]{};
}

\definecolor{silver}{rgb}{0.95,0.95,0.95}

\newtheorem{tdefinitie}{Definitie}[section]
\newenvironment{definitie}[1]%
  {\begin{mdframed}[backgroundcolor=silver,
    topline=false,
    rightline=false,
    leftline=false,
    bottomline=false]\begin{tdefinitie}#1\\\normalfont}%
  {\end{tdefinitie}\end{mdframed}}

\newtheorem{talgo}{Algoritme}[section]
\newenvironment{algoritme}[1]%
  {\begin{mdframed}[backgroundcolor=silver,
    topline=false,
    rightline=false,
    leftline=false,
    bottomline=false]\begin{talgo}#1\\\normalfont}%
  {\end{talgo}\end{mdframed}}
  
\renewcommand{\qedsymbol}{$\blacksquare$}
\newenvironment{bewijs}%
  {\begin{mdframed}[topline=true,
    rightline=true,
    leftline=true,
    bottomline=true]\begin{proof}[Bewijs]\normalfont}%
  {\end{proof}\end{mdframed}}
 
\newenvironment{nfa}%
  {\begin{tikzpicture}[->,
    >=stealth',
    shorten >=1pt,
    auto,
    node distance=2cm,
    semithick]
  \tikzstyle{every state}=[fill=none,draw=black,text=black]
  \tikzstyle{accepting}=[accepting by double]
  \tikzstyle{every initial by arrow}=[initial text=]}%
  {\end{tikzpicture}}
 
\newenvironment{pda}%
  {\begin{tikzpicture}[->,
    >=stealth',
    shorten >=1pt,
    auto,
    node distance=2cm,
    semithick]
  \tikzstyle{every state}=[fill=none,draw=black,text=black]
  \tikzstyle{accepting}=[accepting by double]
  \tikzstyle{every initial by arrow}=[initial text=]}%
  {\end{tikzpicture}}
  
\newenvironment{enumalgo}%
  {\begin{enumerate}[leftmargin=1.5cm,label=Stap \arabic*:]}%
  {\end{enumerate}}
  
\newenvironment{nobreakparagraph}%
  {\par\nobreak\vfil\penalty0\vfilneg\vtop\bgroup}%
  {\par\xdef\tpd{\the\prevdepth}\egroup\prevdepth=\tpd}%
  
\newcommand{\powerset}{\ensuremath{\mathcal{P}}}
\newcommand{\partition}{\ensuremath{\mathfrak{P}}}
\newcommand{\nat}{\ensuremath{\mathbb{N}}}
\newcommand{\dfastate}{\ensuremath{\mathcal{Q}}}
\newcommand{\en}{\ensuremath{\text{ en }}}
\newcommand{\of}{\ensuremath{\text{ of }}}

\begin{document}

\begin{titlepage}
    \newpage
    \thispagestyle{empty}
    \frenchspacing
    \hspace{-0.2cm}
    \includegraphics[height=3.4cm]{assets/sedes}
    \hspace{0.2cm}
    \rule{0.5pt}{3.4cm}
    \hspace{0.2cm}
    \begin{minipage}[b]{8cm}
        \large{Katholieke\newline Universiteit\newline Leuven}\smallskip\newline
        \large{}\smallskip\newline
        \textbf{Department \newline Computerwetenschappen}\smallskip
    \end{minipage}
    \hspace{\stretch{1}}
    \vspace*{3.2cm}\vfill
    \begin{center}
        \begin{minipage}[t]{\textwidth}
            \begin{center}
                \large{\rm{\textbf{\uppercase{Samenvatting}}}}\\
                \large{\rm{Automaten en Berekenbaarheid [G0P84a]}}
            \end{center}
        \end{minipage}
    \end{center}
    \vfill
    \hfill\makebox[3cm][l]{%
        \vbox to 7cm{\vfill\noindent
            {\rm \textbf{Fr\'ed\'eric Hannes}}\\[2mm]
            {\rm Academisch jaar 2015--2016}
        }
    }
\end{titlepage}

\tableofcontents

\newpage

\section{Voorkennis}

\subsection{Griekse letters}

\begin{tabular}{l|l|l}
	Alpha & $\alpha$ & \\
	Beta & $\beta$ & \\
	Delta & $\delta$ & $\Delta$ \\
	Epsilon & $\varepsilon, \epsilon$ & \\
	Gamma & $\gamma$ & $\Gamma$ \\
	Mu & $\mu$ & \\
	Omega & $\omega$ & $\Omega$ \\
	Phi & $\varphi, \phi$ & $\Phi$ \\
	Rho & $\rho$ & \\
	Sigma & $\sigma$ & $\Sigma$ \\
	Tau & $\tau$ & \\
	Theta & $\theta$ & $\Theta$ \\
	Xi & $\xi$ & $\Xi$ \\
\end{tabular}

\subsection{Verzamelingenleer}

\begin{definitie}{Een algebra van een verzameling S}
Een algebra is een verzameling met daarop een aantal inwendige operaties.
\end{definitie}

\begin{definitie}{De unie van twee verzamelingen}
  De unie van twee verzamelingen $S_1$ en $S_2$ is de verzameling die alle elementen bevat die in $S_1$ en/of $S_2$ zitten. \\
  \\ $S_1 \cup S_2 = \{x|x \in S_1 \of x \in S_2\}$
\end{definitie}

\begin{definitie}{De doorsnede van twee verzamelingen}
  De doorsnede van twee verzamelingen $S_1$ en $S_2$ is de verzameling die alle elementen bevat die in zowel $S_1$ als $S_2$ zitten. \\
  \\ $S_1 \cup S_2 = \{x|x \in S_1 \en x \in S_2\}$
\end{definitie}

\begin{definitie}{Het complement van een verzameling}
  Het complement van een verzameling $S$ van alle elementen uit het domein die niet in $S$ zitten. \\
  \\ $\overline{S} = \{x|x \notin S\}$
\end{definitie}

\begin{definitie}{Een machtsverzameling $\powerset(S)$ van een verzameling $S$}
  Een machtsverzameling $\powerset(S)$ is een verzameling van alle deelverzamelingen die gevormd kunnen worden met de elementen in S, inclusief de lege verzameling en de verzameling A.\\ \\ $\powerset(S) = \{A|A \subseteq S\}$
\end{definitie}

\begin{definitie}{Een equivalentierelatie $\sim$ op een verzameling $S$}
  Een equivalentierelatie ($\sim$) is een relatie die de elementen van een verzameling groepeert die allemaal voldoen aan een zekere set criteria. Elke equivalentierelatie is
  \begin{itemize}
  \item reflexief: $\forall x \in S: x \sim x$
  \item symmetrisch: $\forall x, y \in S: x \sim y \Leftrightarrow y \sim x$
  \item transitief: $\forall x, y, z \in S: x \sim y \en y \sim z \Rightarrow x \sim z$
  \end{itemize}
\end{definitie}

\begin{definitie}{Een equivalentieklasse onder $\sim$ op een verzameling $S$}
  Een equivalentieklasse ($x_\sim \subseteq S$) is de verzameling van alle elementen in de verzameling $S$ die onder de gegeven equivalentierelatie gelijkwaardig zijn aan $x$. \\
  \\ $x_\sim = \{y \in S|x \sim y\}$
\end{definitie}

\begin{definitie}{Een partitie $\partition$ van een verzameling $S$}
  Een partitie van een verzameling is een opdeling van de elementen van die verzameling in niet-lege onderling disjuncte delen. Als $\partition$ een partitie is van een verzameling $S$, dan geldt:
  \begin{itemize}
  \item $\emptyset \notin \partition$
  \item $\forall X, Y \in \partition: X \neq Y \Rightarrow X \cap Y = \emptyset$
  \item $\bigcup_{X \in \partition} = S$
  \end{itemize}
\end{definitie}

\begin{definitie}{Een tralie ($L$, $\leq$, $\vee$, $\wedge$)}
  Een tralie is een partieel-geordende verzameling, geordend volgens $\leq$, met een supremum ($\vee$) en een infimum ($\wedge$) voor elk paar elementen. Voor een tralie geldt:
  \begin{itemize}
  \item $\forall x, y \in L: x \leq y \Rightarrow x = x \wedge y$
  \item $\forall x, y \in L: x \leq y \Rightarrow y = x \vee y$
  \end{itemize}
\end{definitie}

\begin{definitie}{Een volledige tralie ($L$, $\leq$, $\vee$, $\wedge$)}
  Een tralie is een volledige tralie indien elke deelverzameling van $L$ een bovengrens en ondergrens heeft die een elementen zijn van $L$. \\
  \\ $\forall X \subseteq L: \bigwedge_X \in L \en \bigvee_X \in L$
\end{definitie}

\newpage\section{Talen}

\begin{definitie}{Een alfabet $\Sigma$}
  Een verzameling van symbolen.
\end{definitie}

\begin{definitie}{Een string $s$ over een alfabet $\Sigma$}
  Een opeenvolging nul, \'e\'en of meerdere elementen uit $\Sigma$. \\
  \\ $s = a_1a_2a_3...a_n$ met $a_i \in \Sigma, n \in \nat$
  \\ $s \in L$
\end{definitie}

\begin{definitie}{De lege string $\epsilon$ over een alfabet $\Sigma$}
  De lege string is een string die bestaat uit nul elementen.
\end{definitie}

\begin{definitie}{Het alfabet $\Sigma_\epsilon$}
  $\Sigma_\epsilon = \Sigma \cup \{\epsilon\}$
\end{definitie}

\begin{definitie}{Een taal $L$ over een alfabet $\Sigma$}
  Een verzameling van strings over $\Sigma$.\\
  \\ $L \subseteq \Sigma^*$
  \\ $L \in L_\Sigma$
\end{definitie}

\begin{definitie}{Concatenatie van twee talen}
  $L_1L_2$ is de concatenatie van twee talen $L_1$ en $L_2$ over hetzelfde alfabet $\Sigma$. \\
  \\ $L_1L_2 = \{xy|x \in L_1, y \in L_2\}$
  \\ $L^n = \{x_1x_2x_3...x_n|x_i \in L, n \in \nat\}$
  \\ $L^0 = \{\epsilon\}$
\end{definitie}

\begin{definitie}{De Kleene ster $L^*$ van een taal $L$}
  $L^* = \bigcup^\infty_{n=0}L^n$
  \\ $L^+ = LL^*$
\end{definitie}

\begin{definitie}{De verzameling $\Sigma^*$ van alle strings over $\Sigma$}
  $\Sigma^*$ is een verzameling met alle strings die gevormd kunnen worden uit het alfabet $\Sigma$. Elke taal over $\Sigma$ is een deelverzameling van $\Sigma^*$. \\
  \\ $\Sigma^* = \{a_1a_2a_3...a_n|a_i \in \Sigma, n \in \nat\}$
\end{definitie}

\begin{definitie}{De verzameling $L_\Sigma$ van alle talen over $\Sigma$}
  $L_\Sigma = \powerset(\Sigma^*)$ is een verzameling met alle talen die gevormd kunnen worden uit het alfabet $\Sigma$. Elke taal over $\Sigma$ is een element van $\powerset(\Sigma^*)$.
\end{definitie}

De verzameling $L_\Sigma$ vormt een algebra met verschillende inwendige operaties, waaronder:
\begin{itemize}
\item Unie: $L_1 \cup L_2 \in L_\Sigma$ met $L_1 \in L_\Sigma, L_2 \in L_\Sigma$
\item Doorsnede: $L_1 \cap L_2 \in L_\Sigma$ met $L_1 \in L_\Sigma, L_2 \in L_\Sigma$
\item Complement: $\overline{L} \in L_\Sigma$ met $L \in L_\Sigma$
\item Concatenatie: $L_1L_2 \in L_\Sigma$ met $L_1 \in L_\Sigma, L_2 \in L_\Sigma$
\item Willekeurig lange concatenatie: $L_1^* \in L_\Sigma$ met $L_1 \in L_\Sigma$
\end{itemize}

\newpage\section{Reguliere talen}

\begin{definitie}{Een reguliere expressie $E$ over een alfabet $\Sigma$}
  $E$ is een reguliere expressie over $\Sigma$ indien $E$ van de vorm is:
  \begin{itemize}
  \item $\epsilon$ (lege string)
  \item $\phi$
  \item $a$ met $a \in \Sigma$
  \item $(E_1E_2)$ met $E_1$ en $E_2$ reguliere expressies over $\Sigma$
  \item $(E_1)^*$ met $E_1$ een reguliere expressie over $\Sigma$
  \item $(E_1|E_2)$ met $E_1$ en $E_2$ reguliere expressies over $\Sigma$
  \end{itemize}
\end{definitie}

\begin{definitie}{Een reguliere expressie $E$ over een alfabet $\Sigma$ bepaalt een taal $L_E$ over $\Sigma$}
  Een taal bepaalt door een reguliere expressie is een reguliere taal.
  \begin{itemize}
  \item Als $E = \epsilon$, dan $L_E = \{\epsilon\}$
  \item Als $E = \phi$, dan $L_E = \emptyset$
  \item Als $E = a$ met $a \in \Sigma$, dan $L_E = \{a\}$
  \item Als $E = (E_1E_2)$, dan $L_E = L_{E_1}L_{E_2}$
  \item Als $E = (E_1)^*$, dan $L_E = L_{E_1}^*$
  \item Als $E = (E_1|E_2)$, dan $L_E = L_{E_1} \cup L_{E_2}$
  \end{itemize}
\end{definitie}

\begin{definitie}{De verzameling $RegLan$ van de reguliere talen}
  $RegLan$ is een verzameling die alle reguliere talen bevat. \\
  \\ $RegLan \subseteq L_\Sigma$
  \\ $x \in RegLan \Rightarrow x \in L_\Sigma$
  \\ $x \in RegLan \en y \in x \Rightarrow y \in \Sigma^*$
\end{definitie}

% TODO (Optional): Bewijs dat RegLan een subalgebra is van L_Sigma

\subsection{NFA}

\begin{definitie}{Een niet-deterministische eindige toestandsautomaat}
  Een niet-deterministische eindige toestandsautomaat (NFA) is een 5-tal ($Q$, $\Sigma$, $\delta$, $q_s$, $F$):
  \begin{itemize}
  \item $Q$ is een eindige verzameling toestanden
  \item $\Sigma$ is een eindig alfabet
  \item \bm{$\delta : Q \times \Sigma_\epsilon \rightarrow \powerset(Q)$}\\$\delta$ is de overgangsfunctie
  \item $q_s \in Q$ is de starttoestand
  \item $F \subseteq Q$ is de verzameling van eindtoestanden
  \end{itemize}
\end{definitie}

Eigenschappen van een NFA:
\begin{itemize}
\item Een NFA is een gerichte graaf.
\item Elke knoop van een NFA heeft een label.
\item Knopen met een dubbele cirkel zijn eindtoestanden.
\item Lussen (bogen van een knoop naar zichzelf) zijn toegestaan in een NFA.
\item Boven van een NFA dragen een label met \'e\'en of meerdere symbolen uit het alfabet, gescheiden door een komma en/of $\epsilon$.
\item Er komt slechts \'e\'en boog aan in de starttoestand, die niet vertrekt vanuit een knoop.
\end{itemize}

% TODO: Uitvoering NFA
% TODO: Algebra van NFA's beschrijven

\begin{definitie}{Een string $s$ wordt aanvaard door een NFA}
  Een string $s$ ($a_1a_2a_3...a_n$ met $a_i \in \Sigma_\epsilon$) wordt aanvaard door een NFA ($Q$, $\Sigma$, $\delta$, $q_s$, $F$) als een rij toestanden $t_1t_2t_3...t_{n+1}$ bestaat zodanig dat:
  \begin{itemize}
  \item $t_1 = q_s$
  \item $t_{i+1} \in \delta(t_i, a_i)$
  \item $t_{n+1} \in F$
  \end{itemize}
\end{definitie}

\begin{definitie}{Een taal bepaald door een NFA $M$}
  Een taal $L$ wordt bepaald door een NFA $M$ als $M$ elke string van $L$ aanvaardt. Deze taal noteren we als $L_M$.
\end{definitie}

\begin{definitie}{Equivalente NFA's}
  Twee NFA's zijn equivalent als ze dezelfde taal bepalen.
\end{definitie}

% TODO: Transitietabel

Elke NFA beschrijft een reguliere taal en elke reguliere taal kan beschreven worden door een NFA ($L_E = L_{NFA_E}$). Figuur \ref{fig:regexnfa} toont de overgangen van de elementaire reguliere expressies naar de bijhorende NFA's.

\begin{figure}[ht]
\label{fig:regexnfa}
\centering
\renewcommand{\arraystretch}{1.5}
\begin{tabular}{>{\centering\arraybackslash}m{2cm}|>{\centering\arraybackslash}m{10cm}}
$E$ & $NFA_E$ \\ \hline
$\epsilon$ & \begin{nfa}
  \node[initial,state,accepting] (A) {};
  \addvmargin{1mm}
\end{nfa} \\
$\phi$ & \begin{nfa}
  \node[initial,state] (A) {};
  \addvmargin{1mm}
\end{nfa} \\
$a \in \Sigma$ & \begin{nfa}
  \node[initial,state]   (A)              {};
  \node[state,accepting] (B) [right of=A] {};
  \path (A) edge [] node {$a$} (B);
  \addvmargin{1mm}
\end{nfa} \\
$(E_1|E_2)$ & \begin{nfa}
  \node[initial,state]   (A)                    []                                     {};
  \node[]                (B)                    [above right=.25cm and .5cm of A,draw] {$NFA_1$};
  \node                  (C)                    [below right=.25cm and .5cm of A,draw] {$NFA_2$};
  \node[state]           (D)                    [right of=B]                           {};
  \node[state]           (E)                    [right of=C]                           {};
  \node[state,accepting] (F) at ($(C)!0.5!(D)$) [right=3.7cm of A]                     {};
  
  \path (A) edge [bend left]  node               {$\epsilon$} (B)
            edge [bend right] node [below left]  {$\epsilon$} (C)
        (B) edge []           node               {}           (D)
        (C) edge []           node               {}           (E)
        (D) edge [bend left]  node               {$\epsilon$} (F)
        (E) edge [bend right] node [below right] {$\epsilon$} (F);
  \addvmargin{1mm}
\end{nfa} \\
$(E_1E_2)$ & \begin{nfa}
  \node[initial]         (A) [draw]            {$NFA_1$};
  \node[state]           (B) [right of=A]      {};
  \node                  (C) [right of=B,draw] {$NFA_2$};
  \node[state,accepting] (D) [right of=C]      {};
  
  \path (A) edge [] node {}           (B)
        (B) edge [] node {$\epsilon$} (C)
        (C) edge [] node {}           (D);
  \addvmargin{1mm}
\end{nfa} \\
$(E_1)^*$ & \begin{nfa}
  \node[initial,state]   (A)                   {};
  \node                  (B) [right of=A,draw] {$NFA_1$};
  \node[state]           (C) [right of=B]      {};
  \node[state,accepting] (D) [right of=C]      {};
  
  \path (A) edge []           node         {$\epsilon$} (B)
            edge [bend right] node [above] {$\epsilon$} (C)
        (B) edge []           node         {$\epsilon$} (C)
        (C) edge []           node         {}           (D)
            edge [bend right] node [above] {$\epsilon$} (A);
  \addvmargin{1mm}
\end{nfa} \\
\end{tabular}
\caption{De overgang van een reguliere expressie naar een NFA.}
\end{figure}

\subsection{GNFA}

Een NFA kan omgezet worden naar reguliere expressie door de NFA eerst om te zetten naar een GNFA, vervolgens een GNFA met twee toestanden, waaruit de reguliere expressie rechstreeks af te lezen is.

\begin{definitie}{Een gegeneraliseerde NFA}
  \label{def:gnfa}
  Een gegeneraliseerde NFA (GNFA) is een NFA met gewijzigde en extra eigenschappen:
  \begin{itemize}
  \item Er is slechts \'e\'en eindtoestand die verschilt van de begintoestand.
  \item Er komen geen bogen aan in de starttoestand (vanuit andere toestanden) en vanuit de starttoestand vertrekt exact \'e\'en boog naar elke andere toestand.
  \item Er is exact \'e\'en boog van elke toestand naar de eindtoestand.
  \item Tussen elk paar toestanden is er exact \'e\'en boog in beide richtingen.
  \item Er is exact \'e\'en boog van elke toestand naar zichzelf.
  \item Het label van elke boog is een reguliere expressie.
  \end{itemize}
\end{definitie}

\begin{algoritme}{Uitvoering van een GNFA}
  \vspace{-5mm}\begin{enumalgo}
  \item Je start een GNFA met een string $s$.
  \item Je kiest een overgang naar een volgende toestand en verwijder de eerste $n$ symbolen die voldoen aan de reguliere expressie op de gevolgde boog. Bij een boog met enkel $\epsilon$ moet er geen symbool verwijderd worden en het is niet toegestaan om een boog te volgen met als label $\phi$.
  \item Als je aankomt in de eindtoestand en de string is leeg, dan heeft de GNFA de string $s$ aanvaard.
  \end{enumalgo}
\end{algoritme}

\begin{algoritme}{Transformatie NFA naar GNFA}
  \vspace{-5mm}\begin{enumalgo}
  \item Generaliseer de NFA:
    \begin{itemize}
    \item Voer een nieuwe starttoestand $q_s$ in.
    \item Voer een nieuwe (unieke) eindtoestand $q_a$ in.
    \item Teken $\epsilon$-boog van $q_s$ naar de oude begintoestand.
    \item Teken $\epsilon$-boog van elke oude eindtoestand naar $q_e$.
    \item Vul de automaat aan met $\phi$-bogen waar er bogen ontbreken. (Zie definitie \ref{def:gnfa})
    \item Neem parallelle gerichte bogen samen met de unie van hun labels.
    \end{itemize}
  \item Reduceer de GNFA:
    \begin{itemize}
    \item Als $Q \setminus \{q_s, q_e\} = \emptyset$, ga naar stap 3.
    \item Voer een reductiestap uit door een willekeurige knoop $q \in Q \setminus \{q_s, q_e\}$ te verwijderen. (Zie voorbeeld \ref{ex:gnfa})
    \item Herhaal stap 2.
    \end{itemize}
  \item Bepaal de reguliere expressie die af te lezen is als label op de boog tussen $q_s$ en $q_e$.
\end{enumalgo}
\end{algoritme}

% TODO: Bewijs dat elke string die aanvaard wordt voor de reductiestap nadien nog steeds aanvaard wordt
% TODO: Bewijs dat elke string die verworpen wordt voor de reductiestap nadien nog steeds verworpen wordt

\subsection{DFA}

\begin{definitie}{Een deterministische eindige toestandsautomaat}
  \label{def:dfa}
  Een deterministische eindige toestandsautomaat (DFA) is een type automaat dat gelijkt op een NFA, waarvoor de volgende extra eigenschappen gelden:
  \begin{itemize}
  \item Er zijn geen bogen met $\epsilon$.
  \item Elk symbool $a \in \Sigma$ mag hoogstens op \'e\'en uitgaande boog staan per knoop.
  \end{itemize}
\end{definitie}

\begin{definitie}{Een complete DFA}
  \label{def:cdfa}
  Een DFA is compleet wanneer de overgangsfunctie $\delta$ totaal is. D.w.z. dat voor elke $a \in \Sigma$ een overgang mogelijk is vanuit elke toestand. Definitie \ref{def:dfa} beschrijft een parti\"ele DFA. Een parti\"ele (niet-complete) DFA kan omgevormd worden tot een complete DFA door een extra toestand $q_t$ toe te voegen waar alle niet-acceptabele strings in vast komen te zitten. We bepalen een nieuwe overgangsfunctie $\delta'$ voor $Q' = Q \cup \{q_t\}$:
  \begin{equation*}
  \forall q \in Q', a \in \Sigma: \delta'(q, a) = \begin{cases}
    \delta(q, a), & \text{als}\ \delta(q, a) \neq \phi\\
    q_t & \text{anders}
  \end{cases}
  \end{equation*}
\end{definitie}

Definitie \ref{def:cdfa} impliceert ook dat voor elke $a \in \Sigma$ een overgang bestaat van $q_t$ naar $q_t$, waarin $t$ staat voor ``trash".

De complete DFA wordt voornamelijk gebruikt in praktijk. Verder in dit document wordt vaak gebruik gemaakt van een DFA met de assumptie dat $\delta$ een totale functie is. Daarbij wordt impliciet gesteld dat voor een niet-complete DFA eerst de constructie van de equivalente complete DFA wordt uitgevoerd.

\begin{algoritme}{Transformatie NFA naar DFA}
  We construeren een DFA ($Q_d$, $\Sigma$, $\delta_d$, $q_{sd}$, $F_d$) uit een NFA ($Q_n$, $\Sigma$, $\delta_n$, $q_{sn}$, $F_n$) zodanig dat $L_{NFA} = L_{DFA}$ als volgt:
  \begin{itemize}
  \item $Q_d = \powerset(Q_n)$: Elke toestand van de DFA is een verzameling van toestanden van de NFA. Onbereikbare toestanden zullen we later verwijderen, waardoor zal gelden: $Q_d \subseteq \powerset(Q_n)$.
  \item \bm{$\delta_d: (\powerset(Q_n) \times \Sigma) \rightarrow \powerset(Q_n)$}\\$\delta_d(\dfastate, a) = eb(\delta_n(\dfastate, a))$ voor $\dfastate \in Q_d$: Voor een symbool $a$ is er vanuit de toestand van de DFA $\dfastate$ een overgang naar de toestand van de DFA die alle toestanden van de NFA bevat die epsilon-bereikbaar zijn vanuit alle $q \in \dfastate$.
  \item $q_{sd} = eb(q_{sn})$: De starttoestand is de verzameling van toestanden van de NFA die epsilon-bereikbaar zijn vanuit de starttoestand van de NFA.
  \item $F_d = \{S|S \in Q_d, S \cap F_n \neq \emptyset\}$: Een eindtoestand van een DFA bevat altijd een eindtoestand van de NFA.
  \end{itemize}
  Om de constructie uit te voeren hebben we de volgende definities ingevoerd:
  \begin{itemize}
  \item \bm{$eb: Q_n \rightarrow \powerset(Q_n)$}\\ $eb(q) = \{x|x \in \{q\} \cup \delta_n(q, \epsilon)\}$ met $q \in Q_n$: De afbeelding die een toestand $q \in Q_n$ afbeeldt op alle epsilon-bereikbare toestanden.
  \item \bm{$eb: Q_d \rightarrow \powerset(Q_n)$}\\ $eb(\dfastate) = \{x|x \in \bigcup_{q \in \dfastate}eb(q)\}$ met $\dfastate \in Q_d$: De afbeelding die een toestand $\dfastate \in Q_d$ afbeeldt op alle toestanden die epsilon-bereikbaar zijn vanuit alle toestanden in $\dfastate$.
  \item \bm{$\delta_n: Q_d \times \Sigma \rightarrow \powerset(Q_n)$}\\ $\delta_n(\dfastate, a) = \{x|x \in \bigcup_{q \in \dfastate}\delta_n(q, a)\}$ met $\dfastate \in Q_d$ en $a \in \Sigma$: De afbeelding die een toestand $\dfastate \in Q_d$ afbeeldt op de verzameling van alle toestanden die voor een gegeven symbool bereikbaar zijn vanuit elke toestand in $\dfastate$.
  \end{itemize}
\end{algoritme}

Elke reguliere taal wordt bepaald door een DFA, omdat elke reguliere taal wordt bepaald door een NFA en elke NFA kan worden omgezet tot een DFA.

\begin{definitie}{De overgangsfunctie $\delta^*$ van een DFA voor strings}
  \label{def:deltaster}
  \bm{$\delta^*: (\powerset(Q) \times \Sigma^*) \rightarrow \powerset(Q)$}\\
  $\delta^*$ is een afbeelding die een string afbeeldt op een verzameling toestanden die vanuit een toestand bereikt kunnen worden met die string. We defini\"eren $\delta^*$ inductief:
  \begin{equation*}
    \delta^*(q, s) = \begin{cases}
      q, & \text{als}\ s=\epsilon \\
      \delta^*(\delta(q, a), w) & \text{als}\ s=aw\ \text{voor}\ a \in \Sigma, w \in \Sigma^*
    \end{cases}
  \end{equation*}
\end{definitie}

\begin{bewijs}
  We bewijzen dat $\delta^*(q, wa) = \delta(\delta^*(q, w), a)$ voor $a \in \Sigma, w \in \Sigma^*_\epsilon$.\\[.2cm]
  Bewijs met inductie op basis van definitie \ref{def:deltaster}:
  \begin{itemize}
  \item Stel $w = \epsilon$, dan geldt:\\$\delta^*(q, \epsilon a) = \delta^*(q, a \epsilon) = \delta^*(\delta(q, a), \epsilon) = \delta(q, a) = \delta(\delta^*(q, \epsilon), a)$
  \item Stel $w \neq \epsilon$ en $w = bw'$ met $b \in \Sigma, w' \in \Sigma^*_\epsilon$, dan geldt:\\$\delta^*(q, wa) = \delta^*(q, bw'a) = \delta^*(\delta(q, b), w'a) = \delta(\delta^*(\delta(q, b), w'), a) = \delta(\delta^*(q, w), a)$
  \end{itemize}
\end{bewijs}

\subsection{Minimale DFA}

\begin{definitie}{f-gelijke toestanden}
  Twee toestanden zijn finaal-gelijk indien
  \vspace{-3mm}\begin{equation*}
  \forall w \in \Sigma^* : \delta^*(p, w) \in F \Leftrightarrow \delta^*(q, w) \in F
  \end{equation*}
\end{definitie}

\begin{definitie}{f-verschillende toestanden}
  Twee toestanden zijn finaal-verschillend indien ze niet f-gelijk zijn.
\end{definitie}

Een minimale DFA wordt opgesteld door f-gelijke toestanden samen te nemen. Om het minimaliseren te vergemakkelijken stellen we een totale functie  $\delta'$ op van $\delta$ waarbij in elke toestand voor elk ontbrekend symbool uit het alfabet een overgang wordt toegevoegd zodanig dat $\delta'(q_i, a)=q_i$.

% TODO: Algoritme minimale DFA
% TODO: Bewijs minimale DFA heeft minimaal aantal toestanden
% TODO: Isomorfe DFA's

\subsection{Myhill-Nerode relaties op $\Sigma^*$}

\begin{definitie}{$\partition_1$ is fijner dan $\partition_2$}
  Als een partie $\partition_1$ fijner is dan een $\partition_2$, dan zit elk element van $\partition_1$ vervat in een element van $\partition_2$. $\{\partition_1$, $\partition_2\}$ vormt een tralie waar geldt dat $\partition_1 \leq \partition_2$.
\end{definitie}

\begin{definitie}{$\partition_1$ is groffer dan $\partition_2$}
  Als een partitie $\partition_1$ groffer is dan een $\partition_2$, dan zit elk element van $\partition_2$ vervat in een element van $\partition_1$. $\{\partition_2$, $\partition_1\}$ vormt een tralie waar geldt dat $\partition_2 \leq \partition_1$.
\end{definitie}

\begin{definitie}{Equivalentierelatie $\sim_L$ over $\Sigma^*$}
  Gegeven een taal $L$ over $\Sigma$, is het mogelijk om $\Sigma^*$ op te delen in de partitie $\{L,\overline{L}\}$. De equivalentierelatie $\sim_L$ over $\Sigma^*$ beschouwt twee elementen van $\Sigma^*$ gelijkwaardig indien ze beiden tot dezelfde deelverzameling van die partitie behoren.
  \begin{equation*}
  \forall x, y \in \Sigma^*: x \sim_L y \Leftrightarrow x,y \in L \of x,y \notin L
  \end{equation*}
\end{definitie}

\begin{definitie}{De deelverzameling $reach(q)$ van $\Sigma^*$ voor een DFA}
  De deelverzameling $reach(q)$ van $\Sigma^*$ voor een DFA ($Q$, $\Sigma$, $\delta$, $q_s$, $F$) bevat alle strings die je vanaf de begintoestand $q_s$ in de toestand $q \in Q$ brengen.
  \begin{equation*}
  reach(q) = \{s \in \Sigma^*|\delta^*(q_s,s)=q\}
  \end{equation*}
\end{definitie}

\begin{definitie}{Equivalentierelatie $\sim_{DFA}$ over $\Sigma^*$}
  Gegeven een $DFA$ met $L_{DFA}$ een taal over $\Sigma$, is het mogelijk om $\Sigma^*$ op te delen in de een partitie $\{reach(q)|q \in Q\}$ zodanig dat
  \begin{itemize}
  \item de partitie eindig is.
  \item $s \in \Sigma^*$ wordt aanvaard door de $DFA$ alss $s \in reach(q)$ voor $q \in F$.
  \item als $s \in reach(q)$ dan $\forall a \in \Sigma : sa \in reach(\delta(q, a))$.
  \end{itemize}
  Of in symbolen:
  \begin{equation*}
  \forall x, y \in \Sigma^*: x \sim_{DFA} y \Leftrightarrow \delta^*(q_s, x) = \delta^*(q_s, y)
  \end{equation*}
\end{definitie}

De equivalentierelatie $\sim_{DFA}$ is een $MN(L_{DFA})$-relatie op $\Sigma^*$.

\begin{definitie}{Een Myhill-Nerode relatie voor een taal $L$ over $\Sigma^*$}
  \label{def:mnl}
  Een Myhill-Nerode relatie voor een taal $L$ ($MN(L)$-relatie) is een equivalentierelatie $\sim_X$ die voldoet aan de volgende eigenschappen:
  \begin{itemize}
  \item $\sim_X$ is rechts congruent\\$\forall x, y \in \Sigma^*, a \in \Sigma: x \sim_X y \Rightarrow xa \sim_X ya$
  \item $\sim_X$ verfijnt $\sim_L$\\$x \sim_X y \Rightarrow x \sim_L y$
  \item $\sim_X$ heeft een eindige index\\Het aantal equivalentieklassen van $\sim_X$ is eindig.
  \end{itemize}
\end{definitie}

\begin{bewijs}
  We bewijzen dat een $MN(L)$-relatie enkel bestaat alss de taal $L$ regulier is.\\[.2cm]
    Het bewijs wordt geleverd door aan te tonen dat het mogelijk is om een $MN(L)$-relatie te construeren uit een DFA en een DFA uit een $MN(L)$-relatie.
  \begin{itemize}
  \item Voor een gegeven DFA $D$ kan men een equivalentierelatie $\sim_D$ construeren waarvoor geldt
  \begin{equation*}
  \forall x, y \in \Sigma^*: x \sim_D y \Leftrightarrow \delta^*(q_s, x) = \delta^*(q_s, y)
  \end{equation*}
  \begin{itemize}
  \item $\sim_D$ is rechts congruent omdat geldt
  \begin{equation*}
  \forall x, y \in \Sigma^*, a \in \Sigma: xa \sim_D ya \Leftrightarrow \delta(\delta^*(q_s, x), a) = \delta(\delta^*(q_s, y), a)
  \end{equation*}
  \item $\sim_D$ verfijnt $\sim_{L_D}$, omdat $\sim_D$ een equivalentieklasse genereert voor alle strings die aanvaard worden door $D$ per eindtoestand en een klasse voor elke string die eindigt in een andere toestand. Daardoor worden de verzamelingen $L$ en $\overline{L}$ opgedeeld in kleinere disjuncte verzamelingen. De verzameling van die equivalentieklassen is een partitie van $\Sigma^*$ die $\sim_{L_D}$ verfijnt.
  \item $\sim_D$ heeft een eindige index omdat het aantal toestanden in $D$ eindig is.
  \end{itemize}
  \item Voor een $MN(L)$-relatie kan men een DFA ($Q$, $\Sigma$, $\delta$, $q_s$, $F$) construeren die $L$ bepaalt, waarvoor geldt
  \begin{itemize}
  \item $Q = \{x_\sim|x \in \Sigma^*\}$\\
  De toestanden van de DFA zijn gelijk aan de equivalentieklassen van de $MN(L)$-relatie, dat aantal equivalentieklassen is eindig.
  \item $\delta(x_\sim, a) = (xa)_\sim$\\
  Volgt uit de rechtse congruentie eigenschap van de equivalentierelatie.
  \item $q_s = \epsilon_\sim$
  Een starttoestand wordt bereikt met $\epsilon$.
  \item $F = \{x_\sim|x \in L\}$\\
  Een eindtoestand wordt bereikt door elke string in $L$. Het aantal eindtoestanden is eindig omdat $F \subseteq Q$.
  \end{itemize}
  Tenslotte bewijzen we met inductie voor een willekeurige string $x \in \Sigma^*$ dat geldt
  \begin{equation*}
  \forall x \in \Sigma^*: x \in L_{DFA} \triangleq \delta^*(\epsilon_\sim,x) \in F \Leftrightarrow x \in L \triangleq x_\sim \in F
  \end{equation*}
  \begin{itemize}
  \item Voor een string $x$ met een lengte $l = 0$:
  \begin{equation*}
  \delta^*(\epsilon_\sim,\epsilon) = \delta(\epsilon_\sim,\epsilon) = \epsilon_\sim
  \end{equation*}
  \item Stel dat de bewering geldt voor een string met willekeurige lengte $l$, dan geldt ze ook voor $xa$ met $a \in \Sigma$:
  \begin{equation*}
  \delta^*(\epsilon_\sim,xa) = \delta(x_\sim,a) = (xa)_\sim
  \end{equation*}
  \end{itemize}
  \end{itemize}
\end{bewijs}

% TODO: MN(L) en isomorfe DFA's
% TODO: Supremum van MN(L)-relaties

Het is mogelijk om met een $MN(L)$-relatie te bewijzen dat een taal niet-regulier is. Je gaat te werk door aan te nemen dat als er een minimale DFA bestaat voor de taal, waarvoor men een relatie $\sim_{DFA}$ definieert. Je neemt aan dat die voldoet aan de eerste twee eigenschappen van de definitie (\ref{def:mnl}) en bewijst vervolgens dat het aantal equivalentieklassen van die relatie oneindig is.

% TODO: Stelling van Myhill en Nerode

\subsection{Pompend lemma voor reguliere talen}

Het pompend lemma kan niet gebruikt worden om te bewijzen dat een taal regulier is. Elke reguliere taal moet voldoen aan de eigenschappen van het pompend lemma, maar er zijn ook niet-reguliere talen die daaraan voldoen. Daardoor kan je met zekerheid zeggen dat een taal niet-regulier is als je ze niet kan pompen, maar je kan niet met zekerheid zeggen dat een taal regulier is. In FO-logica $Regulier(L) \rightarrow KanPompen(L)$ met waarheidstabel:
\begin{center}
\begin{tabular}{r|c|c}
                     & $KanPompen(L)$ & $\neg KanPompen(L)$ \\ \hline
  $Regulier(L)$      & True           & False               \\ \hline
  $\neg Regulier(L)$ & True           & True
\end{tabular}
\end{center}

\begin{definitie}{Het pompend lemma voor reguliere talen}
  Voor een reguliere taal $L$ bestaat een pomplengte $d$ zodanig dat als $s \in L$ en $|s| \geq d$, dan bestaat er minstens \'e\'en verdeling van $s$ in stukken $x$, $y$ en $z$ zodanig dat $s = xyz$ en
  \begin{enumerate}
  \item $\forall i \geq 0: xy^iz \in L$
  \item $|y| > 0$
  \item $|xy| \leq d$
  \end{enumerate}
\end{definitie}

\begin{bewijs}
  We bewijzen de eigenschappen van het pompend lemma. \\[.2cm]
  Voor een DFA die een taal $L$ bepaalt nemen we een pomplengte $d = \#Q + 1$ en een willekeurige string $s = a_1a_2a_3...a_n$ met $n \geq d$. De accepterende sequentie van toestanden ($q_s = q_1$, $q_2$, ..., $q_f$) voor s heeft een lengte strikt groter dan het aantal toestanden in $Q$, dus er zijn bij de eerste $d$ toestanden zeker twee toestanden gelijk omdat er maar $d-1$ toestanden zijn. Stel dat $q_i$ en $q_j$ gelijk zijn met $i < j \leq d$, dan nemen we
  \begin{equation*}
  \begin{cases}
  x = a_1a_2...a_i\\
  y = a_{i+1}a_{i+2}...a_j\\
  z = a_{j+1}a_{j+2}...a_n
  \end{cases}
  \end{equation*}
  \begin{itemize}
  \item Eigenschap 1 geldt omdat $y$ een lus volgt en die $k$ keer herhaald kan worden voor $xy^kz \in L$.
  \item Eigenschap 2 geldt omdat $i$ strikt kleiner is dan $j$, dus bevat $y$ minstens \'e\'en element.
  \item Eigenschap 3 geldt omdat $i < j \leq d$, dus we hebben $xy$ kleiner dan (of gelijk aan) $d$ gekozen.
  \end{itemize}
\end{bewijs}

% TODO: Voorbeeld gebruiken van pompend lemme of aan te tonen dat een taal niet-regulier is
% TODO: Algebra van DFAs

\subsection{Andere automaten}

\begin{definitie}{B\"uchi automaat}
  Een B\"uchi automaat (BA) is een automaat die lijkt op een NFA, maar deze doorlopen oneindige strings. Een BA accepteert een string indien de rij toestanden die gepasseerd wordt oneindig dikwijls een accepterende toestand bevat.
\end{definitie}

\newpage\section{Contextvrije talen}

\begin{definitie}{Contextvrije grammatica}
  Een contextvrije grammatica (CFG) is een 4-tal ($V$, $\Sigma$, $R$, $S$) waarbij
  \begin{itemize}
  \item $V$ een eindige verzameling niet-eindsymbolen is (variabelen of non-terminals)
  \item $\Sigma$ een eindig alfabet van eindsymbolen (terminals), disjunct met $V$
  \item $R$ een eindige verzameling regels (producties) is die bestaan uit een koppel van een variabele en een string $s \in \Sigma^*_\epsilon$, gescheiden door $\rightarrow$
  \item $S \in V$ is de starttoestand
  \end{itemize}
\end{definitie}

\begin{definitie}{Afleiding/constructie m.b.v. een CFG}
  Gegeven een CFG ($V$, $\Sigma$, $R$, $S$) en een string $f$ over $V \cup \Sigma_\epsilon$, dan kan men een string $b$ over $V \cup \Sigma_\epsilon$ afleiden/construeren uit $f$ ($b \Rightarrow^* f$) met de gegeven CFG, indien er er een rij strings $s_0,s_1,s_2,...,s_n$ bestaat zodanig dat
  \begin{itemize}
  \item $s_0 = b$
  \item $s_n = f$
  \item $s_{i+1}$ is gelijk aan $s_i$ (met $i < n$) waarin een variabele $X$ wordt vervangen door de rechterzijde van een productie waarin $X$ aan de linkerzijde voorkomt.
  \end{itemize}
\end{definitie}

\begin{definitie}{Taal $L_{CFG}$ bepaald door een CFG}
  Een taal $L_{CFG}$ bepaald door een CFG ($V$, $\Sigma$, $R$, $S$) is een verzameling van strings over $\Sigma$ die kan afgeleid worden uit het startsymbool $S$.
  \begin{equation*}
  L_{CFG} = \{s \in \Sigma^*|S \Rightarrow^* s\}
  \end{equation*}
\end{definitie}

\begin{definitie}{Contextvrije taal}
  Een taal $L$ is een contextvrije taal (CFL) indien er een CFG bestaat zodanig dat $L = L_{CFG}$.
\end{definitie}

\begin{definitie}{Ambigue strings}
  Een string is ambigu indien er voor die meerdere meest-linkse (of meest-rechtse) afleidingen bestaan.
\end{definitie}

\begin{definitie}{Ambigue CFG's}
  Een CFG is ambigu indien er strings $s \in L_{CFG}$ die ambigue zijn t.o.v. de CFG.
\end{definitie}
Sommige CFL's hebben niet-ambigue CFG, maar dat geldt niet voor elke CFL. Een CFL die geen niet-ambigue CFG heeft is inherent ambigu.

\begin{definitie}{Equivalente CFG's}
  Twee CFG's zijn equivalent indien ze dezelfde taal bepalen.
  \begin{equation*}
  L_{CFG_1} = L_{CFG_2}
  \end{equation*}
\end{definitie}

\begin{definitie}{Chromsky Normaal Vorm}
  Elke CFG heeft een Chromsky Normaal Vorm (CNF) die dezelfde taal genereert en waarbij de producties een beperkt aantal vormen kan aannemen:
  \begin{itemize}
  \item $A \rightarrow BC$ met $\begin{cases}A \in V \\ B,C \in V \setminus \{S\}\end{cases}$
  \item $A \rightarrow \alpha$ met $\begin{cases}A \in V \\ \alpha \in \Sigma\end{cases}$
  \item $S \rightarrow \epsilon$
  \end{itemize}
\end{definitie}

Met de CNF
\begin{itemize}
\item kan je onmiddellijk zien of $L_{CFG}$ de lege string bevat.
\item is elke parse tree (bijna) een volledige binaire boom.
\item heeft de afleiding van een string van lengte $n > 0$ altijd een lengte $2n-1$.
\end{itemize}

% TODO: Bewijs CNF

\begin{definitie}{Greibach Normaal Vorm}
  Elke CFG heeft een Greibach Normaal Vorm (GNF) die dezelfde taal genereert en waarbij de producties een beperkt aantal vormen kan aannemen:
  \begin{itemize}
  \item $A \rightarrow \alpha X$ met $\begin{cases}A \in V \\ \alpha \in \Sigma \\ X = X_1X_2X_3...X_n \text{ met } X_i \in V, n \in \nat\end{cases}$
  \item $S \rightarrow \epsilon$
  \end{itemize}
\end{definitie}

\subsection{PDA}

\begin{definitie}{Push-down automaat}
  Een push-down automaat (PDA) is een 6-tal ($Q$, $\Sigma$, $\Gamma$, $\delta$, $q_s$, $F$) waarbij
  \begin{itemize}
  \item $Q$ is een eindige verzameling toestanden
  \item $\Sigma$ is een eindige inputalfabet
  \item $\Gamma$ is een eindige stapelalfabet
  \item \bm{$\delta: Q \times \Sigma_\epsilon \times \Gamma_\epsilon \rightarrow \powerset(Q \times \Gamma_\epsilon)$}\\$\delta$ is een overgangsfunctie
  \item $q_s$ is de starttoestand
  \item $F \subseteq Q$ is een verzameling eindtoestanden
  \end{itemize}
\end{definitie}

\begin{figure}[!ht]
  \label{fig:pdasstep}
  \centering
  \begin{pda}
  \node[state] (A) []               {};
  \node[state] (B) [right=3cm of A] {};
  
  \path (A) edge [] node {$\alpha,\beta\rightarrow\gamma$} (B);
  \addvmargin{1mm}
  \end{pda}
  \caption{Een overgang van een PDA.}
\end{figure}

\begin{nobreakparagraph}
Een PDA maakt overgangen zoals in figuur \ref{fig:pdasstep},
\begin{itemize}
\item indien $\alpha$ het eerste symbool is van de string
\item en $\beta$ bovenaan de stapel staat,
\item je kan die boog volgen en
\begin{itemize}
\item $\alpha$ van de string verwijderen
\item $\beta$ van de stapel verwijderen
\item $\gamma$ op de stapel plaatsen
\end{itemize}
\end{itemize}
\end{nobreakparagraph}
Indien de PDA niet-deterministisch is, kunnen er meerdere bogen met dezelfde $\alpha,\beta$ combinaties vertrekken uit een toestand.

% TODO: Aanvaarding van een string door een PDA

Alternatieve definities van een PDA:
\begin{itemize}
\item Er mag meer dan \'e\'en symbool op de stapel geplaatst worden tijdens een overgang.
\item Een string wordt geaccepteerd indien de string op is en de stapel leeg is ($F$ heeft geen belang).
\item Een string wordt geaccepteerd indien de string op is, de stapel leeg is en we in een eindtoestand zitten.
\item In een overgang wordt een symbool op de stapel geplaatst maar er wordt geen symbool van de stapel verwijderd of omgekeerd.
\item Er is slechts \'e\'en eindtoestand.
\end{itemize}
Voor een CFL is het mogelijk om met eender welke van deze definities equivalente automaten te maken.

% TODO: Bewijs equivalentie CFG en PDA

\begin{algoritme}{Constructie van een PDA uit een CFG}
  \label{alg:cfgpda}
  We stellen een PDA ($Q$, $\Sigma_{PDA}$, $\Gamma$, $\delta$, $q_s$, $F$) voor een CFG ($V$, $\Sigma_{CFG}$, $R$, $S$) met
  \begin{itemize}
  \item $Q = {q_s, q_h, q_a}$ met $q_h$ een hulptoestand
  \item $\Sigma_{PDA} = \Sigma_{CFG}$
  \item $\Gamma = \{\$\} \cup V \cup \Sigma_{CFG}$
  \item $\delta$ waarvoor geldt:
  \begin{itemize}
  \item Er is 1 boog van $q_s$ naar $q_h$ met label ``$\epsilon,\epsilon\rightarrow\$$".
  \item Er is 1 boog van $q_h$ naar $q_a$ met label ``$\epsilon,\$\rightarrow\epsilon$".
  \item Voor elk symbool $a \in \Sigma_{CFG}$ is er een lus bij $q_h$ met label ``$a,a\rightarrow\epsilon$".
  \item Voor elke productie $(X \rightarrow \gamma) \in R$ is er een lus bij $q_h$ met label ``$\epsilon,X\rightarrow\gamma$". (Vervang de linkerkant van een productie bovenaan de stapel met de rechterkant van de productie.)
  \end{itemize}
  \item $q_s$ de starttoestand
  \item $F = \{q_a\}$ met $q_a$ de eindtoestand
  \end{itemize}
\end{algoritme}

Om een meest-linkse afleiding te bekomen moeten we in algoritme \ref{alg:cfgpda} de symbolen van een string $\gamma \in \Sigma^*_\epsilon$ in omgekeerde volgorde op de stapel plaatsen.

% TODO: Een PDA accepteert een L_CFG
% TODO: Constructie CFG uit PDA

\subsection{Pompend lemma voor contextvrije talen}

\section{Voorbeelden}

\subsection{Transitietabel NFA}
\label{ex:nfatbl}

\begin{nfa}
  \node[initial,state]   (S)                    {$S$};
  \node[state]           (A) [below right of=S] {$A$};
  \node[state,accepting] (E) [above right of=A] {$E$};
  
  \path (S) edge []           node         {$a$}          (E)
            edge []           node         {$a,\epsilon$} (A)
        (A) edge [loop below] node         {$b$}          (A)
            edge [bend right] node         {$a$}          (E)
        (E) edge [loop right] node         {$a,b$}        (E)
            edge [bend right] node [above] {$\epsilon$}   (A);
  \addvmargin{1mm}
\end{nfa}

De bijhorende transitietabel voor de gegeven NFA:

\begin{center}
\begin{tabular}{r|r|r}
  $Q$ & $\Sigma_\epsilon$ & $\powerset(Q)$ \\ \hline
  $S$ & $a$ & $\{A,E\}$ \\
  $S$ & $\epsilon$ & $\{A\}$ \\
  $A$ & $a$ & $\{E\}$ \\
  $A$ & $b$ & $\{A\}$ \\
  $E$ & $a$ & $\{E\}$ \\
  $E$ & $b$ & $\{E\}$ \\
  $E$ & $\epsilon$ & $\{A\}$ \\ \hline
  $S$ & $b$ & $\emptyset$ \\
  $A$ & $\epsilon$ & $\emptyset$ \\
\end{tabular}
\end{center}

\subsection{NFA naar RE}
\label{ex:gnfa}

\begin{nfa}
  \node[initial,state]   (S)                    {$S$};
  \node[state]           (A) [below right of=S] {$A$};
  \node[state,accepting] (E) [above right of=A] {$E$};
  
  \path (S) edge []           node         {$a$}          (E)
            edge []           node         {$a,\epsilon$} (A)
        (A) edge [loop below] node         {$b$}          (A)
            edge [bend right] node         {$a$}          (E)
        (E) edge [loop right] node         {$a,b$}        (E)
            edge [bend right] node [above] {$\epsilon$}   (A);
  \addvmargin{1mm}
\end{nfa}

\paragraph{Stap 1} Generaliseren

\begin{nfa}
  \node[initial,state]   (SN)                    {$q_s$};
  \node[state]           (S)  [right of=SN]      {$S$};
  \node[state]           (A)  [below right of=S] {$A$};
  \node[state]           (E)  [above right of=A] {$E$};
  \node[state,accepting] (EN) [right of=E]       {$q_e$};
  
  \path (SN) edge []           node         {$\epsilon$}   (S)
        (S)  edge []           node         {$a$}          (E)
             edge []           node         {$a|\epsilon$} (A)
        (A)  edge [loop below] node         {$b$}          (A)
             edge [bend right] node         {$a$}          (E)
        (E)  edge [loop above] node         {$a|b$}        (E)
             edge [bend right] node [above] {$\epsilon$}   (A)
             edge []           node         {$\epsilon$}   (EN);
  \addvmargin{1mm}
\end{nfa}

\paragraph{Stap 2} Reductie (met vereenvoudiging)

Reductie 1:\\
\begin{nfa}
  \node[initial,state]   (SN)                    {$q_s$};
  \node[state]           (S)  [right of=SN]      {$S$};
  \node[state]           (A)  [below right of=S] {$A$};
  \node[state]           (E)  [above right of=A] {$E$};
  \node[state,accepting] (EN) [right of=E]       {$q_e$};
  
  \path (SN) edge []           node {$\epsilon$}   (S)
        (S)  edge []           node {$a$}          (E)
             edge []           node {$a|\epsilon$} (A)
        (A)  edge [loop below] node {$b$}          (A)
             edge []           node {$a$}          (E)
        (E)  edge [loop above] node {$a|b|b^*a$}   (E)
             edge []           node {$\epsilon$}   (EN);
  \addvmargin{1mm}
\end{nfa}

\begin{nfa}
  \node[initial,state]   (SN)                  {$q_s$};
  \node[state]           (S)  [right of=SN]    {$S$};
  \node[state]           (E)  [right=2cm of S] {$E$};
  \node[state,accepting] (EN) [right of=E]     {$q_e$};
  
  \path (SN) edge []           node {$\epsilon$}         (S)
        (S)  edge [bend right] node {$(a|\epsilon)b^*a$} (E)
             edge [bend left]  node {$a$}                (E)
        (E)  edge [loop above] node {$b|b^*a$}           (E)
             edge []           node {$\epsilon$}         (EN);
  \addvmargin{1mm}
\end{nfa}

\begin{nfa}
  \node[initial,state]   (SN)                  {$q_s$};
  \node[state]           (S)  [right of=SN]    {$S$};
  \node[state]           (E)  [right=3cm of S] {$E$};
  \node[state,accepting] (EN) [right of=E]     {$q_e$};
  
  \path (SN) edge []           node {$\epsilon$}         (S)
        (S)  edge []           node {$(a|\epsilon)b^*a$} (E)
        (E)  edge [loop above] node {$b|b^*a$}              (E)
             edge []           node {$\epsilon$}         (EN);
  \addvmargin{1mm}
\end{nfa}

Reductie 2:\\
\begin{nfa}
  \node[initial,state]   (SN)                  {$q_s$};
  \node[state]           (S)  [right of=SN]    {$S$};
  \node[state,accepting] (EN) [right=4cm of S] {$q_e$};
  
  \path (SN) edge [] node {$\epsilon$}                   (S)
        (S)  edge [] node {$(a|\epsilon)b^*a(b|b^*a)^*$} (EN);
  \addvmargin{1mm}
\end{nfa}

\begin{nfa}
  \node[initial,state]   (SN)                  {$q_s$};
  \node[state]           (S)  [right of=SN]    {$S$};
  \node[state,accepting] (EN) [right=4cm of S] {$q_e$};
  
  \path (SN) edge [] node {$\epsilon$}                   (S)
        (S)  edge [] node {$(a|\epsilon)b^*a(a|b)^*$} (EN);
  \addvmargin{1mm}
\end{nfa}

Reductie 3:\\
\begin{nfa}
  \node[initial,state]   (SN)                   {$q_s$};
  \node[state,accepting] (EN) [right=4cm of SN] {$q_e$};
  
  \path (SN) edge [] node {$(a|\epsilon)b^*a(a|b)^*$} (EN);
  \addvmargin{1mm}
\end{nfa}

\begin{nfa}
  \node[initial,state]   (SN)                   {$q_s$};
  \node[state,accepting] (EN) [right=4cm of SN] {$q_e$};
  
  \path (SN) edge [] node {$b^*a(a|b)^*$} (EN);
  \addvmargin{1mm}
\end{nfa}

\paragraph{Stap 3} De RE voor de NFA is ``$b^*a(a|b)^*$".

\subsection{NFA naar DFA}
\label{ex:nfadfa}

\subsubsection{Voorbeeld 1}

\begin{nfa}
  \node[initial,state]   (S)                    {$S$};
  \node[state]           (A) [below right of=S] {$A$};
  \node[state,accepting] (E) [above right of=A] {$E$};
  
  \path (S) edge []           node         {$a$}          (E)
            edge []           node         {$a,\epsilon$} (A)
        (A) edge [loop below] node         {$b$}          (A)
            edge [bend right] node         {$a$}          (E)
        (E) edge [loop right] node         {$a,b$}        (E)
            edge [bend right] node [above] {$\epsilon$}   (A);
  \addvmargin{1mm}
\end{nfa}

We construeren een tabel met alle toestanden van $Q_d = \powerset(Q_n)$ en de epsilon-bereikbare functie:

\begin{center}
\begin{tabular}{r|r}
  $q \in Q_d$ & $eb(q)$ \\ \hline
  $\{\}$ & $\{\}$ \\
  $\{S\}$ & $\{S,A\}$ \\
  $\{A\}$ & $\{A\}$ \\
  $\{E\}$ & $\{A,E\}$ \\
  $\{S,A\}$ & $\{S,A\}$ \\
  $\{S,E\}$ & $\{S,A,E\}$ \\
  $\{A,E\}$ & $\{A,E\}$ \\
  $\{S,A,E\}$ & $\{S,A,E\}$ \\
\end{tabular}
\end{center}

We stellen nu de verzamelingen voor $\delta_n(Q_d,a)$ en $\delta_n(Q_d,b)$ op:

\begin{center}
\begin{tabular}{r|r|r|r}
  $q \in Q_d$ & $eb(q)$ & $\delta_n(q,a)$ & $\delta_n(q,b)$ \\ \hline
  $\{\}$ & $\{\}$ & $\{\}$ & $\{\}$ \\
  $\{S\}$ & $\{S,A\}$ & $\{A,E\}$ & $\{\}$ \\
  $\{A\}$ & $\{A\}$ & $\{E\}$ & $\{A\}$ \\
  $\{E\}$ & $\{A,E\}$ & $\{E\}$ & $\{E\}$ \\
  $\{S,A\}$ & $\{S,A\}$ & $\{A,E\}$ & $\{A\}$ \\
  $\{S,E\}$ & $\{S,A,E\}$ & $\{A,E\}$ & $\{E\}$ \\
  $\{A,E\}$ & $\{A,E\}$ & $\{E\}$ & $\{A,E\}$ \\
  $\{S,A,E\}$ & $\{S,A,E\}$ & $\{A,E\}$ & $\{A,E\}$ \\
\end{tabular}
\end{center}

Tenslotte stellen nu de verzamelingen voor $\delta_d(Q_d,a)=eb(\delta_n(Q_d,a))$ en $\delta_d(Q_d,b)=eb(\delta_n(Q_d,b))$ op:

\begin{center}
\begin{tabular}{r|r|r|r|r|r}
  $q \in Q_d$ & $eb(q)$ & $\delta_n(q,a)$ & $\delta_n(q,b)$ & $\delta_d(q,a)$ & $\delta_d(q,b)$ \\ \hline
  $\{\}$ & $\{\}$ & $\{\}$ & $\{\}$ & $\{\}$ & $\{\}$ \\
  $\{S\}$ & $\{S,A\}$ & $\{A,E\}$ & $\{\}$ & $\{A,E\}$ & $\{\}$ \\
  $\{A\}$ & $\{A\}$ & $\{E\}$ & $\{A\}$ & $\{A,E\}$ & $\{A\}$ \\
  $\{E\}$ & $\{A,E\}$ & $\{E\}$ & $\{E\}$ & $\{A,E\}$ & $\{A,E\}$ \\
  $\{S,A\}$ & $\{S,A\}$ & $\{A,E\}$ & $\{A\}$ & $\{A,E\}$ & $\{A\}$ \\
  $\{S,E\}$ & $\{S,A,E\}$ & $\{A,E\}$ & $\{E\}$ & $\{A,E\}$ & $\{A,E\}$ \\
  $\{A,E\}$ & $\{A,E\}$ & $\{E\}$ & $\{A,E\}$ & $\{A,E\}$ & $\{A,E\}$ \\
  $\{S,A,E\}$ & $\{S,A,E\}$ & $\{A,E\}$ & $\{A,E\}$ & $\{A,E\}$ & $\{A,E\}$ \\
\end{tabular}
\end{center}

De resulterende DFA:

\begin{nfa}
  \node[initial,state]   (S)                    {$\{S,A\}$};
  \node[state]           (A) [below right of=S] {$\{A\}$};
  \node[state,accepting] (E) [above right of=A] {$\{A,E\}$};
  
  \path (S) edge []           node {$a$}   (E)
            edge []           node {$b$}   (A)
        (A) edge [loop below] node {$b$}   (A)
            edge []           node {$a$}   (E)
        (E) edge [loop right] node {$a,b$} (E);
  \addvmargin{1mm}
\end{nfa}

\subsubsection{Voorbeeld 2}

\begin{nfa}
  \node[initial,state]   (q1)               {$q_1$};
  \node[state,accepting] (q2) [right of=q1] {$q_2$};
  
  \path (q1) edge [loop above] node {$a,b$} (q1)
             edge []           node {$b$}   (q2);
  \addvmargin{1mm}
\end{nfa}

We construeren een tabel met alle toestanden van $Q_d = \powerset(Q_n)$ en de epsilon-bereikbare functie:

\begin{center}
\begin{tabular}{r|r}
  $q \in Q_d$ & $eb(q)$ \\ \hline
  $\{\}$ & $\{\}$ \\
  $\{q_1\}$ & $\{q_1\}$ \\
  $\{q_2\}$ & $\{q_2\}$ \\
  $\{q_1,q_2\}$ & $\{q_1,q_2\}$ \\
\end{tabular}
\end{center}

We stellen nu de verzamelingen voor $\delta_n(Q_d,a)$ en $\delta_n(Q_d,b)$ op:

\begin{center}
\begin{tabular}{r|r|r|r}
  $q \in Q_d$ & $eb(q)$ & $\delta_n(q,a)$ & $\delta_n(q,b)$ \\ \hline
  $\{\}$ & $\{\}$ & $\{\}$ & $\{\}$ \\
  $\{q_1\}$ & $\{q_1\}$ & $\{q_1\}$ & $\{q_1,q_2\}$ \\
  $\{q_2\}$ & $\{q_2\}$ & $\{\}$ & $\{\}$ \\
  $\{q_1,q_2\}$ & $\{q_1,q_2\}$ & $\{q_1\}$ & $\{q_1,q_2\}$ \\
\end{tabular}
\end{center}

Tenslotte stellen nu de verzamelingen voor $\delta_d(Q_d,a)=eb(\delta_n(Q_d,a))$ en $\delta_d(Q_d,b)=eb(\delta_n(Q_d,b))$ op:

\begin{center}
\begin{tabular}{r|r|r|r|r|r}
  $q \in Q_d$ & $eb(q)$ & $\delta_n(q,a)$ & $\delta_n(q,b)$ & $\delta_d(q,a)$ & $\delta_d(q,b)$ \\ \hline
  $\{\}$ & $\{\}$ & $\{\}$ & $\{\}$ & $\{\}$ & $\{\}$ \\
  $\{q_1\}$ & $\{q_1\}$ & $\{q_1\}$ & $\{q_1,q_2\}$ & $\{q_1\}$ & $\{q_1,q_2\}$ \\
  $\{q_2\}$ & $\{q_2\}$ & $\{\}$ & $\{\}$ & $\{\}$ & $\{\}$ \\
  $\{q_1,q_2\}$ & $\{q_1,q_2\}$ & $\{q_1\}$ & $\{q_1,q_2\}$ & $\{q_1\}$ & $\{q_1,q_2\}$ \\
\end{tabular}
\end{center}

De resulterende DFA:

\begin{nfa}
  \node[initial,state]   (q1)                     {$\{q_1\}$};
  \node[state,accepting] (q2) [right=1.5cm of q1] {$\{q_1,q_2\}$};
  
  \path (q1) edge [loop above] node         {$a$} (q1)
             edge [bend left]  node [above] {$b$} (q2)
        (q2) edge [loop above] node         {$b$} (q2)
             edge [bend left]  node [below] {$a$} (q1);
  \addvmargin{1mm}
\end{nfa}

\subsection{Parti\"ele DFA naar complete DFA}
\label{ex:dfatotal}

\textit{In dit voorbeeld laten we de bogen van $q \in F$ naar $q_t$ weg om de leesbaarheid te bevorderen. Het is echter belangrijk deze niet te vergeten om een echte complete DFA te vormen!}

\begin{nfa}
  \node[initial,state]   (q1)               {$q_1$};
  \node[state]           (q3) [right of=q1] {$q_3$};
  \node[state]           (q2) [above of=q3] {$q_2$};
  \node[state,accepting] (q4) [right of=q3] {$q_4$};
  \node[state]           (q5) [below of=q3] {$q_5$};
  
  \path (q1) edge [] node {$a$} (q2)
             edge [] node {$b$} (q3)
             edge [] node {$c$} (q5)
        (q2) edge [] node {$c$} (q4)
        (q3) edge [] node {$c$} (q4)
        (q5) edge [] node {$d$} (q4);
  \addvmargin{1mm}
\end{nfa}

Voor de gegeven DFA doen we de omzetting $\delta_{partieel} \longrightarrow \delta_{totaal}$. We voegen een ``trash'' toestand toe waar we alle ontbrekende bogen naar afleiden.

\begin{nfa}
  \node[initial,state]   (q1)               {$q_1$};
  \node[state]           (q3) [right of=q1] {$q_3$};
  \node[state]           (q2) [above of=q3] {$q_2$};
  \node[state,accepting] (q4) [right of=q3] {$q_4$};
  \node[state]           (q5) [below of=q3] {$q_5$};
  \node[state,dashed]    (qt) [left of=q1]  {$q_t$};
  
  \path (q1) edge []                  node              {$a$}       (q2)
             edge []                  node              {$b$}       (q3)
             edge []                  node              {$c$}       (q5)
             edge [dashed,bend right] node [above]      {$d$}       (qt)
        (q2) edge []                  node              {$c$}       (q4)
             edge [dashed,bend right] node [above left] {$a,b,d$}   (qt)
        (q3) edge []                  node              {$c$}       (q4)
             edge [dashed,bend left]  node              {$a,b,d$}   (qt)
        (q5) edge []                  node              {$d$}       (q4)
             edge [dashed,bend left]  node              {$a,b,c$}   (qt)
        (qt) edge [dashed,loop left]  node              {$a,b,c,d$} (qt);
  \addvmargin{1mm}
\end{nfa}

De resulterende automaat is nog steeds een geldige DFA, maar heeft nu de toestand $q_t$ van waaruit het niet mogelijk is om de eindtoestand $q_4$ te bereiken.

\subsection{DFA f-gelijk}
\label{ex:dfafeq}

We willen de f-gelijke toestanden van een DFA bepalen waarvan we eerst de parti\"ele overgangsfunctie totaal hebben gemaakt. We nemen het resultaat van het voorbeeld in sectie \ref{ex:dfatotal}.

We beginnen met de initiatie van het algoritme om de f-gelijke toestanden te bepalen en stellen de initi\"ele graaf $V$.

\begin{ugraph}
  \node[state]        (q1)                     {$q_1$};
  \node[state]        (q2) [above right of=q1] {$q_2$};
  \node[state]        (q3) [right of=q2]       {$q_3$};
  \node[state]        (q4) [below right of=q3] {$q_4$};
  \node[state]        (q5) [below left of=q4]  {$q_5$};
  \node[state,dashed] (qt) [left of=q5]        {$q_t$};
  
  \path (q1) edge []       node {} (q4)
        (q2) edge []       node {} (q4)
        (q3) edge []       node {} (q4)
        (q5) edge []       node {} (q4)
        (qt) edge [dashed] node {} (q4);
  \addvmargin{1mm}
\end{ugraph}

We kiezen twee willekeurige knopen die nog ``onbeslist'' zijn, of m.a.w. die nog niet verbonden zijn in $V$. We nemen $q_1$ en $q_5$ en vinden dat $(\delta(q_1, d), \delta(q_5, d)) = (q_t, q_4) \in V$. We voegen de boog $(q_1, q_5)$ toe aan $V$.

\begin{ugraph}
  \node[state]        (q1)                     {$q_1$};
  \node[state]        (q2) [above right of=q1] {$q_2$};
  \node[state]        (q3) [right of=q2]       {$q_3$};
  \node[state]        (q4) [below right of=q3] {$q_4$};
  \node[state]        (q5) [below left of=q4]  {$q_5$};
  \node[state,dashed] (qt) [left of=q5]        {$q_t$};
  
  \path (q1) edge []       node {} (q4)
             edge []       node {} (q5)
        (q2) edge []       node {} (q4)
        (q3) edge []       node {} (q4)
        (q5) edge []       node {} (q4)
        (qt) edge [dashed] node {} (q4);
  \addvmargin{1mm}
\end{ugraph}

We kiezen opnieuw twee willekeurige knopen die nog ``onbeslist'' zijn, of m.a.w. die nog niet verbonden zijn in $V$. We nemen $q_1$ en $q_2$ en vinden dat $(\delta(q_1, c), \delta(q_5, c)) = (q_5, q_4) \in V$. We voegen de boog $(q_1, q_2)$ toe aan $V$.

\begin{ugraph}
  \node[state]        (q1)                     {$q_1$};
  \node[state]        (q2) [above right of=q1] {$q_2$};
  \node[state]        (q3) [right of=q2]       {$q_3$};
  \node[state]        (q4) [below right of=q3] {$q_4$};
  \node[state]        (q5) [below left of=q4]  {$q_5$};
  \node[state,dashed] (qt) [left of=q5]        {$q_t$};
  
  \path (q1) edge []       node {} (q2)
             edge []       node {} (q4)
             edge []       node {} (q5)
        (q2) edge []       node {} (q4)
        (q3) edge []       node {} (q4)
        (q5) edge []       node {} (q4)
        (qt) edge [dashed] node {} (q4);
  \addvmargin{1mm}
\end{ugraph}

We herhalen de iteratiestap tot we er geen enkel paar knopen meer voldoet aan de voorwaarden en we de volgende graaf $V$ bekomen:

\begin{ugraph}
  \node[state]        (q1)                     {$q_1$};
  \node[state]        (q2) [above right of=q1] {$q_2$};
  \node[state]        (q3) [right of=q2]       {$q_3$};
  \node[state]        (q4) [below right of=q3] {$q_4$};
  \node[state]        (q5) [below left of=q4]  {$q_5$};
  \node[state,dashed] (qt) [left of=q5]        {$q_t$};
  
  \path (q1) edge []       node {} (q2)
             edge []       node {} (q3)
             edge []       node {} (q4)
             edge []       node {} (q5)
        (q2) edge []       node {} (q4)
             edge []       node {} (q5)
        (q3) edge []       node {} (q4)
             edge []       node {} (q5)
        (q5) edge []       node {} (q4)
        (qt) edge [dashed] node {} (q1)
             edge [dashed] node {} (q2)
             edge [dashed] node {} (q3)
             edge [dashed] node {} (q4)
             edge [dashed] node {} (q5);
  \addvmargin{1mm}
\end{ugraph}

Vervolgens bepalen we de complementgraaf $G$ van $V$:

\begin{ugraph}
  \node[state]        (q1)                     {$q_1$};
  \node[state]        (q2) [above right of=q1] {$q_2$};
  \node[state]        (q3) [right of=q2]       {$q_3$};
  \node[state]        (q4) [below right of=q3] {$q_4$};
  \node[state]        (q5) [below left of=q4]  {$q_5$};
  \node[state,dashed] (qt) [left of=q5]        {$q_t$};
  
  \path (q2) edge [] node {} (q3);
  \addvmargin{1mm}
\end{ugraph}

De complementgraaf $G$ bevat een boog tussen $q_2$ en $q_3$, welke een component vormen en dus onderling f-gelijk zijn. Alle andere componenten bevatten slechts een enkele toestand.

\subsection{DFA f-gelijk (partities verfijnen)}
\label{ex:dfafeqpart}

We willen de f-gelijke toestanden van een DFA bepalen waarvan we eerst de parti\"ele overgangsfunctie totaal hebben gemaakt. We nemen het resultaat van het voorbeeld in sectie \ref{ex:dfatotal}.

De eerste partitie wordt opgesteld als een partitie die de toestanden in de eind- en niet-eindtoestanden verdeelt:

\begin{equation*}
P_0 = \{\{q_1,q_2,q_3,q_5,q_t\},\{q_4\}\}
\end{equation*}

Nu voeren we iteratiestappen uit. We partitioneren elke deelverzameling van de partities zodanig dat voor elk symbool, er vanuit beiden toestanden een overgang is naar dezelfde deelverzameling van de partitie.

\begin{equation*}
P_1 = \{\{q_1\},\{q_2,q_3\},\{q_5\},\{q_t\},\{q_4\}\}
\end{equation*}

Voor de gegeven DFA loopt het partitioneren af na een enkele stap. Het is niet mogelijk om de partitie nog verder te verfijnen. Elke deelverzameling van de partitie bevat nu de f-gelijke toestanden van de DFA.

\subsection{DFA minimaliseren}
\label{ex:dfamin}

We zullen een DFA minimaliseren, zoals die gegeven is aan het begin van sectie \ref{ex:dfatotal}. We moeten geen knopen verwijderen bij deze DFA. Het compleet maken van de DFA wordt uitgevoerd in sectie \ref{ex:dfatotal} en in sectie \ref{ex:dfafeq} (of \ref{ex:dfafeqpart}) worden de f-gelijke toestanden bepaalt. We stellen de volgende DFA op:

\begin{nfa}
  \node[initial,state]   (q1)                     {$q_1$};
  \node[state]           (q2) [above right of=q1] {$q_{23}$};
  \node[state,accepting] (q4) [below right of=q2] {$q_4$};
  \node[state]           (q5) [below right of=q1] {$q_5$};
  \node[state,dashed]    (qt) [left of=q1]        {$q_t$};
  
  \path (q1) edge []                  node              {$a,b$}     (q2)
             edge []                  node              {$c$}       (q5)
             edge [dashed,bend right] node [above]      {$d$}       (qt)
        (q2) edge []                  node              {$c$}       (q4)
             edge [dashed,bend right] node [above left] {$a,b,d$}   (qt)
        (q5) edge []                  node              {$d$}       (q4)
             edge [dashed,bend left]  node              {$a,b,c$}   (qt)
        (qt) edge [dashed,loop left]  node              {$a,b,c,d$} (qt);
  \addvmargin{1mm}
\end{nfa}

Tenslotte verwijderen we nog de ``trash''-toestand (of hulptoestand) die we eerder hebben toegevoegd om een totale overgangsfunctie te bekomen. De resulterende DFA is minimaal:

\begin{nfa}
  \node[initial,state]   (q1)                     {$q_1$};
  \node[state]           (q2) [above right of=q1] {$q_{23}$};
  \node[state,accepting] (q4) [below right of=q2] {$q_4$};
  \node[state]           (q5) [below right of=q1] {$q_5$};
  
  \path (q1) edge [] node {$a,b$} (q2)
             edge [] node {$c$}   (q5)
        (q2) edge [] node {$c$}   (q4)
        (q5) edge [] node {$d$}   (q4);
  \addvmargin{1mm}
\end{nfa}

\subsection{Taal met NFA kleiner dan minimale DFA}
\label{ex:nfaldfamin}

Een voorbeeld van een taal waarvoor een NFA bestaat die kleiner is dan DFA$_{min}$ is de taal over $\Sigma = \{a,b\}$ bepaald door een reguliere expressie gegeven als
\begin{equation*}
  (a|b)^*a(a|b)(a|b)
\end{equation*}

Een NFA die dezelfde taal bepaalt is:

\begin{nfa}
  \node[initial,state]   (S)              {$S$};
  \node[state]           (A) [right of=S] {$A$};
  \node[state]           (B) [right of=A] {$B$};
  \node[state,accepting] (E) [right of=B] {$E$};
  
  \path (S) edge [loop above] node {$a,b$} (S)
            edge []           node {$a$}   (A)
        (A) edge []           node {$a,b$} (B)
        (B) edge []           node {$a,b$} (E);
  \addvmargin{1mm}
\end{nfa}

We stellen een verkorte tabel op beginnende bij toestand $S$. Er is geen enkele boog met $\epsilon$, dus $\forall a \in \Sigma: \delta_n(q,a)=\delta_d(q,a)$. Elke toestand van $Q_d$ krijgt een alias om de leesbaarheid te bevorderen.

\begin{center}
\begin{tabular}{r|r|r|r}
  $q \in Q_d$ & $eb(q)$ & $\delta_n(q,a)=\delta_d(q,a)$ & $\delta_n(q,b)=\delta_d(q,b)$ \\ \hline
  $1=\{S\}$ & $\{S\}$ & $\{S,A\}$ & $\{S\}$\\
  $2=\{S,A\}$ & $\{S,A\}$ & $\{S,A,B\}$ & $\{S,B\}$\\
  $3=\{S,A,B\}$ & $\{S,A,B\}$ & $\{S,A,B,E\}$ & $\{S,B,E\}$\\
  $4=\{S,B\}$ & $\{S,B\}$ & $\{S,A,E\}$ & $\{S,E\}$\\
  $5=\{S,A,B,E\}$ & $\{S,A,B,E\}$ & $\{S,A,B,E\}$ & $\{S,B,E\}$\\
  $6=\{S,B,E\}$ & $\{S,B,E\}$ & $\{S,A,E\}$ & $\{S,E\}$\\
  $7=\{S,A,E\}$ & $\{S,A,E\}$ & $\{S,A,B\}$ & $\{S,B\}$\\
  $8=\{S,E\}$ & $\{S,E\}$ & $\{S,A\}$ & $\{S\}$\\
\end{tabular}
\end{center}

We kunnen vervolgende de bijhorende DFA construeren:

\begin{nfa}
  \node[initial,state]   (1)                    {$1$};
  \node[state]           (2) [above left of=1]  {$2$};
  \node[state]           (4) [above right of=2] {$4$};
  \node[state,accepting] (7) [right of=4]       {$7$};
  \node[state]           (3) [right of=7]       {$3$};
  \node[state,accepting] (5) [right of=3]       {$5$};
  \node[state,accepting] (6) [below right of=5] {$6$};
  \node[state,accepting] (8) [below left of=6]  {$8$};
  
  \path (1) edge [loop below]    node {$b$} (1)
            edge []              node {$a$} (2)
        (2) edge []              node {$b$} (4)
            edge [bend left=90]  node {$a$} (3)
        (4) edge [bend left=40]  node {$a$} (7)
            edge [bend right=10] node {$b$} (8)
        (7) edge []              node {$a$} (3)
            edge []              node {$b$} (4)
        (3) edge []              node {$a$} (5)
            edge [bend right=10] node {$b$} (6)
        (5) edge [loop above]    node {$a$} (5)
            edge []              node {$b$} (6)
        (6) edge [bend left=15]  node {$a$} (7)
            edge []              node {$b$} (8)
        (8) edge []              node {$a$} (2)
            edge []              node {$b$} (1);
  \addvmargin{1mm}
\end{nfa}

We proberen om de gevonden DFA nog verder te minimaliseren door de f-verschillende toestanden te bepalen met het standaard algoritme. We bekomen de volgende graaf:

\begin{ugraph}
  \node[state] (1)              {$1$};
  \node[state] (2) [above right of=1] {$2$};
  \node[state] (4) [right of=2] {$4$};
  \node[state] (7) [below right of=4] {$7$};
  \node[state] (3) [below of=7] {$3$};
  \node[state] (5) [below left of=3] {$5$};
  \node[state] (6) [left of=5]  {$6$};
  \node[state] (8) [above left of=6]  {$8$};
  
  % INIT
  \path (1) edge [] node {} (7)
        (2) edge [] node {} (7)
        (4) edge [] node {} (7)
        (3) edge [] node {} (7)
        (1) edge [] node {} (5)
        (2) edge [] node {} (5)
        (4) edge [] node {} (5)
        (3) edge [] node {} (5)
        (1) edge [] node {} (6)
        (2) edge [] node {} (6)
        (4) edge [] node {} (6)
        (3) edge [] node {} (6)
        (1) edge [] node {} (8)
        (2) edge [] node {} (8)
        (4) edge [] node {} (8)
        (3) edge [] node {} (8);
  % REPEAT
  \path (1) edge [] node {} (3)
        (1) edge [] node {} (4)
        (2) edge [] node {} (3)
        (2) edge [] node {} (4)
        (5) edge [] node {} (7)
        (5) edge [] node {} (8)
        (6) edge [] node {} (7)
        (6) edge [] node {} (8)
        (7) edge [] node {} (8)
        (1) edge [] node {} (2)
        (5) edge [] node {} (6);
  \addvmargin{1mm}
\end{ugraph}

Vervolgens bepalen we de complementgraaf:

\begin{ugraph}
  \node[state] (1)              {$1$};
  \node[state] (2) [above right of=1] {$2$};
  \node[state] (4) [right of=2] {$4$};
  \node[state] (7) [below right of=4] {$7$};
  \node[state] (3) [below of=7] {$3$};
  \node[state] (5) [below left of=3] {$5$};
  \node[state] (6) [left of=5]  {$6$};
  \node[state] (8) [above left of=6]  {$8$};
  \addvmargin{1mm}
\end{ugraph}

De complementgraaf bevat geen enkele component die uit meer dan \'e'en toestand bestaat. We concluderen dat de DFA reeds minimaal is. We zien nu ook dat de NFA minder toestanden heeft dan de equivalente minimale DFA.

\subsection{Taal is niet regulier met pompend lemma}
\label{ex:plre}

We tonen aan dat een taal $L = \{a^nb^n|n \in \nat\}$ over $\Sigma = \{a,b\}$ niet regulier is.

Stel dat er voor $L$ een pomplengte $d$ bestaat, dan beschouwen we de string $s=a^db^d$. We nemen een willekeurige opdeling van $s=xyz$ met $|y| > 0$. Er zijn drie mogelijkheden:

\begin{itemize}
\item $y$ bevat enkel $a$'s: $xyyz$ bevat dan meer $a$'s dan $b$'s, dus $xyyz \notin L$ of $xz$ bevat dan meer $b$'s dan $a$'s, dus $xz \notin L$
\item $y$ bevat enkel $b$'s: $xyyz$ bevat dan meer $b$'s dan $a$'s, dus $xyyz \notin L$ of $xz$ bevat dan meer $a$'s dan $b$'s, dus $xz \notin L$
\item $y$ is van de vorm $a^ib^j$ met $i,j \in \nat_0$: $xyyz \notin L$ want het aantal $a$'s en $b$'s blijft enkel gelijk indien $i=j$ en zelfs dan wordt de volgorde niet meer gerespecteerd, er zullen $a$'s en $b$'s door elkaar staan.
\end{itemize}

\subsection{CFG}
\label{ex:cfg}

We stellen een CFG op voor de contextvrije taal $L = \{a^nb^n|n \in \nat\}$ over $\Sigma = \{a,b\}$.

\begin{itemize}
\item $S \rightarrow aSb$
\item $S \rightarrow \epsilon$
\end{itemize}

\subsection{CNF}
\label{ex:cnf}

We stellen een CFG in CNF op voor de contextvrije taal $L = \{a^nb^n|n \in \nat\}$ over $\Sigma = \{a,b\}$.

\begin{itemize}
\item $A \rightarrow a$
\item $B \rightarrow b$
\item $C \rightarrow XB$
\item $X \rightarrow AC$
\item $X \rightarrow AB$
\item $S \rightarrow X$
\item $S \rightarrow \epsilon$
\end{itemize}

\subsection{PDA}
\label{ex:pda}

We stellen een PDA op voor de contextvrije taal $L = \{a^nb^n|n \in \nat\}$ over $\Sigma = \{a,b\}$.

\begin{pda}
  \node[initial,state]   (q1)               {$q_1$};
  \node[state]           (q2) [right of=q1] {$q_2$};
  \node[state]           (q3) [right of=q2] {$q_3$};
  \node[state,accepting] (q4) [right of=q3] {$q_4$};
  
  \path (q1) edge []           node {$\pdat{\epsilon}{\epsilon}{\$}$}       (q2)
        (q2) edge [loop above] node {$\pdat{a}{\epsilon}{a}$}               (q2)
             edge []           node {$\pdat{\epsilon}{\epsilon}{\epsilon}$} (q3)
        (q3) edge [loop above] node {$\pdat{b}{a}{\epsilon}$}               (q3)
             edge []           node {$\pdat{\epsilon}{\$}{\epsilon}$}       (q4);
  \addvmargin{1mm}
\end{pda}

Deze PDA kan niet met minder toestanden voorgesteld worden. Indien $q_3$ een eindtoestand zou zijn (en $q_2$ de starttoestand), dan zou het aantal $a$'s en $b$'s niet meer gelijk zijn, omdat de PDA dan kan stoppen voor dat het even veel $b$'s als $a$'s is tegengekomen. De strings $aaa$ of $aaab$ zouden dan onder andere geldig zijn, zo lang er minder dan of evenveel $b$'s als $a$'s zijn.

\subsection{CFG naar PDA}
\label{ex:cfgpda}

\begin{itemize}
\item $E \rightarrow a$
\item $E \rightarrow E+E$
\item $E \rightarrow E*E$
\end{itemize}

We converteren de gegeven CFG met startsymbool $E$ naar een PDA:

\begin{pda}
  \node[initial,state]   (qs)                   {$q_s$};
  \node[state]           (x)  [below=2cm of qs] {$x$};
  \node[state,accepting] (qf) [below=2cm of x]  {$q_f$};
  
  \path (qs) edge []                            node [pos=0.2] {$\pdat{\epsilon}{\epsilon}{E\$}$} (x)
        (x)  edge []                            node [pos=0.8] {$\pdat{\epsilon}{\$}{\epsilon}$}  (qf)
             edge [out=120,in=150,distance=2cm] node [left]    {$\pdat{\epsilon}{E}{a}$}          (x)
             edge [out=165,in=195,distance=2cm] node [left]    {$\pdat{\epsilon}{E}{E*E}$}        (x)
             edge [out=210,in=240,distance=2cm] node [left]    {$\pdat{\epsilon}{E}{E+E}$}        (x)
             edge [out=60,in=30,distance=2cm]   node [right]   {$\pdat{a}{a}{\epsilon}$}          (x)
             edge [out=15,in=345,distance=2cm]  node [right]   {$\pdat{*}{*}{\epsilon}$}          (x)
             edge [out=330,in=300,distance=2cm] node [right]   {$\pdat{+}{+}{\epsilon}$}          (x);
  \addvmargin{1mm}
\end{pda}

\subsection{Taal is niet contextvrij met pompend lemma}
\label{ex:plcfg}

Stel dat er voor een taal $L = \{a^nb^nc^n|n\geq\}$ over $\Sigma = \{a,b,c\}$ een pomplengte $p$ bestaat en stel dat we voor een string $s = a^pb^pc^p$ met $|s| \geq p$ een opdeling $s=uvxyz$ maken met $|vy| > 0$. Dan zijn er slechts twee mogelijkheden voor die opdeling:
\begin{itemize}
\item $v$ is van de vorm $\alpha^i$ en $y$ van de vorm $\beta^j$ waarbij $\alpha,\beta \in \Sigma$ met $k+l>0$: $uv^2xy^2z \notin L$ kan niet bestaan uit een gelijk aantal $a$'s, $b$'s en $c$'s.
\item $v$ en/of $y$ bestaan uit meer dan \'e\'en symbool uit $\Sigma$: $uv^2xy^2z \notin L$ of meer specifiek $v^2$ en/of $y^2$ bevatten symbolen die niet in de juiste volgorde staan.
\end{itemize}

We hebben aangetoond dat de taal $L$ niet contextvrij is.
\section{Examenvragen}

\subsection{GNFA}

\subsubsection{Vraag 1 (Bijvraag)}

\begin{center}
\renewcommand{\arraystretch}{1.5}
\begin{tabular}{>{\centering\arraybackslash}m{5cm}>{\centering\arraybackslash}m{1cm} >{\centering\arraybackslash}m{5cm}}
\begin{nfa}
  \node[state] (A)                     {$A$};
  \node[state] (X)  [below right of=A] {$X$};
  \node[state] (B)  [above right of=X] {$B$};
  
  \path (A) edge [bend left]  node {$E_4$} (B)
            edge [bend right] node {$E_1$} (X)
        (X) edge [loop below] node {$E_2$} (E)
            edge [bend right] node {$E_3$} (B);
  \addvmargin{1mm}
\end{nfa} & $\longrightarrow$ & \begin{nfa}
  \node[state] (A)                   {$A$};
  \node[state] (B)  [right=3cm of A] {$B$};
  
  \path (A) edge []  node {$E_4|E_1(E_2)^*E_3$} (B);
  \addvmargin{1mm}
\end{nfa}
\end{tabular}
\end{center}

\textit{Stel dat het linkse deel toestand B niet bevat, de boog met $E_3$ van $X$ naar $A$ gaat en de boog met $E_4$ van $A$ naar zichzelf gaat. Kan dit omgevormd worden op methode gelijkaardig aan die beschreven is in de cursus?}

\begin{center}
\begin{nfa}
  \node[state] (A)              {$A$};
  \node[state] (X) [right of=A] {$X$};
  
  \path (A) edge [bend left]  node {$E_1$} (X)
            edge [loop left]  node {$E_4$} (A)
        (X) edge [bend left]  node {$E_3$} (A)
            edge [loop right] node {$E_2$} (X);
  \addvmargin{1mm}
\end{nfa}
\end{center}

Ja, we elimineren eerst het pad van $X$ naar $A$. We voeren een nieuwe toestand $Y$ in en gebruiken de standaard eleminatie om het pad van $Y$ naar $X$ door $A$ te behandelen:

\begin{center}
\renewcommand{\arraystretch}{1.5}
\begin{tabular}{>{\centering\arraybackslash}m{5cm}>{\centering\arraybackslash}m{1cm} >{\centering\arraybackslash}m{5cm}}
\begin{nfa}
  \node[state] (A)                   {$A$};
  \node[state] (X) [above right=.25cm and 2cm of A] {$X$};
  \node[state] (Y) [below right=.25cm and 2cm of A] {$Y$};
  
  \path (A) edge [bend left]  node {$E_1$}      (X)
            edge [loop left]  node {$E_4$}      (A)
        (X) edge [bend left]  node {$\epsilon$} (Y)
            edge [loop right] node {$E_2$}      (X)
        (Y) edge [bend left]  node {$E_3$}      (A)
            edge [bend left]  node {$\epsilon$} (X);
  \addvmargin{1mm}
\end{nfa} & $\longrightarrow$ & \begin{nfa}
  \node[state] (X) []                {$X$};
  \node[state] (Y) [below=.75cm of X] {$Y$};
  
  \path (X) edge [bend left]  node {$\epsilon$}               (Y)
            edge [loop right] node {$E_2$}                    (X)
        (Y) edge [bend left]  node {$\epsilon|E_3(E_4)^*E_1$} (X);
  \addvmargin{1mm}
\end{nfa}
\end{tabular}
\end{center}

Nu we de boog van $X$ naar $A$ ge\"elemineerd hebben, kunnen we de resterende bewerkingen uitvoeren om de lussen van de knopen op de boog tussen $A$ en $X$ te plaatsen:

\begin{center}
\renewcommand{\arraystretch}{1.5}
\begin{tabular}{>{\centering\arraybackslash}m{5cm}>{\centering\arraybackslash}m{1cm} >{\centering\arraybackslash}m{5cm}}
\begin{nfa}
  \node[state] (A)              {$A$};
  \node[state] (X) [right of=A] {$X$};
  
  \path (A) edge []           node {$E_1$} (X)
            edge [loop left]  node {$E_4$} (A)
        (X) edge [loop above] node {$E_2|E_3(E_4)^*E_1$} (X);
  \addvmargin{1mm}
\end{nfa} & $\longrightarrow$ & \begin{nfa}
  \node[state] (A)              {$A$};
  \node[state] (X) [right=4cm of A] {$X$};
  
  \path (A) edge [] node {$(E_4)^*E_1(E_2|E_3(E_4)^*E_1)^*$} (X);
  \addvmargin{1mm}
\end{nfa}
\end{tabular}
\end{center}

\subsection{$MN(L)$-relaties}

\subsubsection{Vraag 1}

\textit{Geef de definitie van een Myhil-Nerode relatie over een taal $L$, of zoals we noteren een $MN(L)$-relatie.}

Een Myhill-Nerode relatie voor een taal $L$ ($MN(L)$-relatie) is een equivalentierelatie $\sim_X$ die voldoet aan de volgende eigenschappen:
  \begin{itemize}
  \item $\sim_X$ is rechts congruent\\$\forall x, y \in \Sigma^*, a \in \Sigma: x \sim_X y \Rightarrow xa \sim_X ya$
  \item $\sim_X$ verfijnt $\sim_L$\\$x \sim_X y \Rightarrow x \sim_L y$
  \item $\sim_X$ heeft een eindige index\\Het aantal equivalentieklassen van $\sim_X$ is eindig.
  \end{itemize}

\textit{Bewijs vervolgens dat een $MN(L)$-relatie bestaat als en slechts als $L$ regulier is.}

  Het bewijs wordt geleverd door aan te tonen dat het mogelijk is om een $MN(L)$-relatie te construeren uit een DFA en een DFA uit een $MN(L)$-relatie.
  \begin{itemize}
  \item Voor een gegeven DFA $D$ kan men een equivalentierelatie $\sim_D$ construeren waarvoor geldt
  \begin{equation*}
  \forall x, y \in \Sigma^*: x \sim_D y \Leftrightarrow \delta^*(q_s, x) = \delta^*(q_s, y)
  \end{equation*}
  \begin{itemize}
  \item $\sim_D$ is rechts congruent omdat geldt
  \begin{equation*}
  \forall x, y \in \Sigma^*, a \in \Sigma: xa \sim_D ya \Leftrightarrow \delta(\delta^*(q_s, x), a) = \delta(\delta^*(q_s, y), a)
  \end{equation*}
  \item $\sim_D$ verfijnt $\sim_{L_D}$, omdat $\sim_D$ een equivalentieklasse genereert voor alle strings die aanvaard worden door $D$ per eindtoestand en een klasse voor elke string die eindigt in een andere toestand. Daardoor worden de verzamelingen $L$ en $\overline{L}$ opgedeeld in kleinere disjuncte verzamelingen. De verzameling van die equivalentieklassen is een partitie van $\Sigma^*$ die $\sim_{L_D}$ verfijnt.
  \item $\sim_D$ heeft een eindige index omdat het aantal toestanden in $D$ eindig is.
  \end{itemize}
  \item Voor een $MN(L)$-relatie kan men een DFA ($Q$, $\Sigma$, $\delta$, $q_s$, $F$) construeren die $L$ bepaalt, waarvoor geldt
  \begin{itemize}
  \item $Q = \{x_\sim|x \in \Sigma^*\}$\\
  De toestanden van de DFA zijn gelijk aan de equivalentieklassen van de $MN(L)$-relatie, dat aantal equivalentieklassen is eindig.
  \item $\delta(x_\sim, a) = (xa)_\sim$\\
  Volgt uit de rechtse congruentie eigenschap van de equivalentierelatie.
  \item $q_s = \epsilon_\sim$
  Een starttoestand wordt bereikt met $\epsilon$.
  \item $F = \{x_\sim|x \in L\}$\\
  Een eindtoestand wordt bereikt door elke string in $L$. Het aantal eindtoestanden is eindig omdat $F \subseteq Q$.
  \end{itemize}
  Tenslotte bewijzen we met inductie voor een willekeurige string $x \in \Sigma^*$ dat geldt
  \begin{equation*}
  \forall x \in \Sigma^*: x \in L_{DFA} \triangleq \delta^*(\epsilon_\sim,x) \in F \Leftrightarrow x \in L \triangleq x_\sim \in F
  \end{equation*}
  \begin{itemize}
  \item Voor een string $x$ met een lengte $l = 0$:
  \begin{equation*}
  \delta^*(\epsilon_\sim,\epsilon) = \delta(\epsilon_\sim,\epsilon) = \epsilon_\sim
  \end{equation*}
  \item Stel dat de bewering geldt voor een string met willekeurige lengte $l$, dan geldt ze ook voor $xa$ met $a \in \Sigma$:
  \begin{equation*}
  \delta^*(\epsilon_\sim,xa) = \delta(x_\sim,a) = (xa)_\sim
  \end{equation*}
  \end{itemize}
  \end{itemize}

\textit{Bestaat er voor een taal $L$ soms meer dan één $MN(L)$-relatie?}

Een taal $L$ waarvoor een $MN(L)$-relatie bestaat moet regulier zijn. Een mogelijke $MN(L)$-relatie voor een reguliere taal is $\sim_{DFA}$. Die equivalentierelatie wordt bepaald door alle strings die vanuit de begintoestand van de $DFA$ een willekeurige toestand bereiken te groeperen per toestand.
\begin{equation*}
\forall x, y \in \Sigma^*: x \sim_{DFA} y \Leftrightarrow \delta^*(q, x) = \delta^*(q, y)
\end{equation*}
We kunnen voor de taal $L$ een equivalente DFA opstellen die dezelfde taal bepaalt, maar niet isomorf is. De equivalentierelatie $\sim_{DFA}$ voor die DFA zal $\Sigma^*$ anders partitioneren. Bijgevolg bestaat er meer dan \'e\'en $MN(L)$-relatie voor een taal $L$.

\end{document}
